% =========================================================================
% Capítulo N: <Nombre del capítulo>
% Tema YY: <Nombre del tema>
% Problema ZZ: <Título del problema>
% Origen: <Ayudantía X, 2026 / Interrogación / Propio>
% =========================================================================
% PLANILLA DE PROBLEMA
%  1. Copia este archivo a Contenidos/CN_<Nombre>/Ejercicios/Problemas/
%     con el nombre CN_EYY_PZZ_titulo_en_snake_case.tex (PZZ = siguiente libre).
%  2. Cáblalo en Ejercicios/CN_EYY_problemas.tex:
%       \input{Contenidos/CN_<Nombre>/Ejercicios/Problemas/CN_EYY_PZZ_titulo}
%       \clearpage
%  3. Borra los bloques que no uses (patos, resumen, etc.).
%
%  Interruptores (main.tex): \mostrarSolucionfalse, \mostrarTipsfalse,
%  \mostrarOjofalse, \mostrarTrucosfalse, \mostrarResumenfalse.
% =========================================================================
\subsubsection{Título del problema}

% --- ENUNCIADO --------------------------------------------------------------
Describe el sistema y las variables, haciendo referencia a la Figura
\ref{fig:cN_eYY_pZZ_titulo}. Determine:

\begin{enumerate}
    \item Primera pregunta.
    \item Segunda pregunta.
\end{enumerate}

% --- FIGURA PRINCIPAL (TikZ nativo, nada de \includegraphics) ---------------
\begin{figure}[!ht]
    \centering
    \begin{tikzpicture}[>=Latex, scale=1.5]
        % --- CONFIGURACIÓN DE PARÁMETROS ---
        \def\R{2.0}   % radio [u]
        \def\ang{35}  % ángulo de dibujo [°]

        % --- SOPORTES Y ELEMENTOS FIJOS ---
        \fill[pattern=north east lines] (-0.5,-0.2) rectangle (0.5,0);
        \draw[thick] (-0.5,0) -- (0.5,0);

        % --- GEOMETRÍA PRINCIPAL ---
        \coordinate (O) at (0,0);
        \coordinate (P) at (\ang:\R);
        \draw[thick] (O) -- (P);

        % --- SISTEMAS DE COORDENADAS ---

        % --- PUNTOS Y COORDENADAS ---
        \filldraw (O) circle (1.5pt) node[below left] {$O$};
        \filldraw (P) circle (1.5pt) node[above right] {$P$};

        % --- VECTORES Y FUERZAS ---

        % --- ANOTACIONES Y ETIQUETAS ---
        % Ángulos: arco + nodo (las comillas de la librería quotes no
        % funcionan con babel en español).
        \draw[dashed] (O) -- (1.2,0);
        \draw[->] (0.8,0) arc (0:\ang:0.8) node[midway, right] {$\theta$};
    \end{tikzpicture}
    \caption{Descripción breve de la figura (45--70 caracteres).}
    \label{fig:cN_eYY_pZZ_titulo}
\end{figure}

% --- SOLUCIÓN (se oculta con \mostrarSolucionfalse) --------------------------
\begin{solucion}{Título del problema}

\textbf{Marcos de referencia y Sistemas de Coordenadas}
\begin{itemize}
    \item $\mathcal{N}$: suelo (inercial) $\longrightarrow \{\vec{E}_x;\vec{E}_y;\vec{E}_z\}$
    \item $\mathcal{A}$: <cuerpo> $\longrightarrow \{\vec{e}_r;\vec{e}_\theta;\vec{E}_z\}$
\end{itemize}

% Figura de apoyo (opcional): sin figure, sin caption, sin label.
% \begin{center}
% \begin{tikzpicture}[scale=0.8] \Polaresrtz{\theta}{E}{e}{azul}{myred} \end{tikzpicture}
% \end{center}

% Transformaciones entre bases: equation* + split
\begin{equation*}
\begin{split}
    \vec{e}_r &= \cos\theta\,\vec{E}_x + \sin\theta\,\vec{E}_y\\
    \vec{e}_\theta &= -\sin\theta\,\vec{E}_x + \cos\theta\,\vec{E}_y
\end{split}
\end{equation*}

$$\velang{N}{A} = \dot\theta\,\vec{E}_z$$

\subsubsection*{Cinemática de la partícula}
\textbf{Posición}\\
$$\vec{r}_{P/O} = r\,\vec{e}_r$$

\textbf{Velocidad}\\
$$\marc{\vec{v}_P}{N} = \der{\vec{r}_{P/O}}{N} = \der{\vec{r}_{P/O}}{A}{} + \velang{N}{A}\times\vec{r}_{P/O}$$

\textbf{Aceleración}\\
$$\marc{\vec{a}_P}{N} = \der{\marc{\vec{v}_P}{N}{}}{N}{}$$

% OJO: \marc y \der se tragan el token siguiente; ciérralas con {} al final
% de un grupo o antes de $$.

% --- Cinética (capítulos C2 en adelante) ---
% \subsubsection*{Cinética}
% \textbf{Diagrama de cuerpo libre}\\
% \textbf{Segunda ley de Newton}\\

% --- Resultado final destacado (en vez de \boxed) ---
\respuesta{$\displaystyle \marc{\vec{a}_P}{N} = (\ddot r - r\dot\theta^2)\,\vec{e}_r + (r\ddot\theta + 2\dot r\dot\theta)\,\vec{e}_\theta$}

% --- PATOS (cada uno con su interruptor) -------------------------------------
% Accesorios extra opcionales entre corchetes: squareglasses, glasses,
% sunglasses, shovel, tie=<color>, book, think=<texto>, megaphone, ...
\patotip[squareglasses]{Consejo general para abordar el problema.}

\patoojo{Error típico o trampa del enunciado.}

\patotruco{Atajo que ahorra álgebra.}

\begin{resumen}[Resumen del método]
Lo que hay que llevarse de este problema.
\end{resumen}

\end{solucion}
